%----------------------------------------------------------------------------------------
% NOTATIE EN FIGUREN VAN DE CURSUS LIMIETEN EN AFGELEIDEN
% Staat in de preambule en niet in een hoofdstuk: ook de lijsten met oplossingen gebruiken dit.
%----------------------------------------------------------------------------------------
% Bewaarde grafiek, geplaatst door maak-pdf.py: bestand en pagina in figuren-cache/
\newcommand{\grafiekbestand}[2]{\includegraphics[page=#2,trim=3bp 3bp 3bp 3bp]{figuren-cache/#1}}

\providecommand{\R}{\mathbb{R}}
\newcommand{\dom}{\operatorname{dom}}
\newcommand{\ber}{\operatorname{ber}}
\newcommand{\inv}[1]{#1^{-1}}

% Zoals stappen, maar met een bredere kolom voor de berekening: voor lange formules
\newenvironment{stappenbreed}{%
  \par\renewcommand{\arraystretch}{1.25}%
  \noindent\begin{tabular}{@{}>{$\displaystyle}p{9.8cm}<{$}@{\hspace{4mm}}>{\RaggedRight\small\itshape\arraybackslash}p{4.9cm}@{}}%
}{\end{tabular}\par}

% Zoals stappen, met eigen kolombreedtes: \begin{stappenvrij}{6.2cm}{3.8cm}. Voor naast een figuur.
\newenvironment{stappenvrij}[2]{%
  \par\renewcommand{\arraystretch}{1.25}%
  \noindent\begin{tabular}{@{}>{$\displaystyle}p{#1}<{$}@{\hspace{3mm}}>{\RaggedRight\small\itshape\arraybackslash}p{#2}@{}}%
}{\end{tabular}\par}

% Tekst of berekening links, figuur rechts: \naast[breedte rechts]{links}{rechts}
% De breedte rechts komt uit de figuur zelf; het optionele argument wordt niet gebruikt.
\newsavebox{\naastbox}
\newcommand{\naast}[3][]{%
  \par\noindent\sbox{\naastbox}{#3}%
  \begin{minipage}[t]{\dimexpr\linewidth-\wd\naastbox-3mm\relax}\vspace{0pt}#2\end{minipage}\hfill
  \begin{minipage}[t]{\wd\naastbox}\vspace{0pt}\usebox{\naastbox}\end{minipage}\par\addvspace{1mm}}
% Een rij figuren, gecentreerd
\newcommand{\figuren}[1]{\par\addvspace{1mm}{\centering #1\par}\addvspace{1.5mm}}
% Opgave met een grafiek rechts: \begin{metgrafiek}[breedte]{tikz} opgave en deelvragen \end{metgrafiek}
% De solution komt na \end{metgrafiek}, zoals bij metprent.
\newsavebox{\grafiekbox}
\newenvironment{metgrafiek}[2][]{%
  \sbox{\grafiekbox}{#2}%
  \begin{minipage}[t]{\dimexpr\linewidth-\wd\grafiekbox-3mm\relax}%
}{%
  \end{minipage}\hfill
  \begin{minipage}[t]{\wd\grafiekbox}\vspace{-2mm}\usebox{\grafiekbox}\end{minipage}\par}

% Rooster met assen: \rooster{xmin}{xstap}{xmax}{ymin}{ystap}{ymax}, met getallen op de assen (gehele waarden).
% \roosterleeg doet hetzelfde zonder getallen, voor eigen labels.
\newcommand{\roosterleeg}[6]{%
  \draw[black!14,very thin,xstep=#2,ystep=#5] (#1,#4) grid (#3,#6);
  \draw[->,grijs] (#1,0) -- ({#3+.4*#2},0) node[right,text=black] {$x$};
  \draw[->,grijs] (0,#4) -- (0,{#6+.4*#5}) node[above,text=black] {$y$};}
\newcommand{\rooster}[6]{%
  \roosterleeg{#1}{#2}{#3}{#4}{#5}{#6}
  \pgfmathtruncatemacro{\rtweex}{#1+#2}\pgfmathtruncatemacro{\rtweey}{#4+#5}
  \foreach \t in {#1,\rtweex,...,#3}{\ifnum\t=0\else\node[below,inner sep=1.5pt,text=black] at (\t,0) {$\t$};\fi}
  \foreach \t in {#4,\rtweey,...,#6}{\ifnum\t=0\else\node[left,inner sep=1.5pt,text=black] at (0,\t) {$\t$};\fi}}

% In de functievoorschriften is \x de veranderlijke; schrijf (\x)^2.
\newcommand{\kromme}[4][]{\draw[accent,thick,#1] plot[domain=#2:#3,samples=90,smooth] (\x,{#4});}
% Een punt op een grafiek: \vol{(x,y)} hoort erbij, \open{(x,y)} niet
\newcommand{\vol}[1]{\fill[accent] #1 circle (2.2pt);}
\newcommand{\open}[1]{\filldraw[draw=accent,thick,fill=white] #1 circle (2.2pt);}
% Tekentabel: \tekentabel{kolommen}{x-rij}{f-rij}, bv. \tekentabel{cccccc}{x & & -2 & & 3 & }{f(x) & - & 0 & + & 0 & -}
\newcommand{\tekentabel}[3]{{\renewcommand{\arraystretch}{1.2}$\begin{array}{c|#1}#2\\\hline #3\end{array}$}}

\providecommand{\e}{\mathrm{e}}
% Rechter- en linkerlimiet, met de ongelijkheid onder de pijl (zoals in Integralen): \limr{x}{0} en \liml{x}{1}
\newcommand{\limr}[2]{\lim\limits_{\substack{#1 \to #2 \\ >}}}
\newcommand{\liml}[2]{\lim\limits_{\substack{#1 \to #2 \\ <}}}
\newcommand{\Z}{\mathbb{Z}}
